% !TeX program = lualatex
% !TeX root = pred1.tex
% ============================================================
% Zajednički preamble za sva predavanja RWA kolegija
% Uključiti s: % ============================================================
% Zajednički preamble za sva predavanja RWA kolegija
% Uključiti s: % ============================================================
% Zajednički preamble za sva predavanja RWA kolegija
% Uključiti s: \input{../common/preamble.tex}
% ============================================================

\documentclass[aspectratio=149,11pt]{beamer}
\usetheme{CambridgeUS}

\usepackage{fontspec}              % radi samo sa XeLaTeX/LuaLaTeX
\setsansfont{Latin Modern Sans}    % sans-serif font
\setmonofont{Latin Modern Mono}    % monospace font za kod
\usepackage[croatian]{babel}
\usepackage{csquotes}              % \enquote{} → hrvatski navodnici „..."
\usepackage{lmodern}
\usepackage{booktabs}
\usepackage{listings}
\usepackage{xcolor}
\usepackage{hyperref}

% ---------- Autor i logo ----------
\author{mag.ing. Josip Torić (SIT UNIZD)}
\titlegraphic{\includegraphics[height=2.5cm]{../common/SIT_logo_tekst1.pdf}}

% ---------- Listings ----------
\definecolor{codebg}{RGB}{248,248,248}
\definecolor{codeframeclr}{RGB}{220,220,220}
\lstdefinelanguage{bash}{
  morekeywords={sudo,apt,brew,pip,pipx,uv,winget,python,git,flask,source,deactivate,code,gunicorn,docker,compose,heroku,render,railway,node,npm,pnpm,yarn,uvicorn,alembic,pytest,ruff,black,precommit},
  sensitive=false,
  morecomment=[l]{\#}
}
\lstdefinelanguage{python}{
  morekeywords={from,import,def,class,return,if,elif,else,for,in,while,try,except,raise,with,as,not,and,or,is,lambda,yield,True,False,None,await,async},
  sensitive=true,
  morecomment=[l]{\#},
  morestring=[b]',
  morestring=[b]"
}
\lstdefinelanguage{sql}{
  morekeywords={SELECT,FROM,WHERE,INSERT,INTO,VALUES,UPDATE,SET,DELETE,CREATE,TABLE,PRIMARY,KEY,FOREIGN,REFERENCES,UNIQUE,NOT,NULL,INDEX,ON,ALTER,DROP,CONSTRAINT},
  sensitive=false,
  morecomment=[l]{--}
}
\lstdefinelanguage{json}{
  morecomment=[l]{//},
  morestring=[b]",
  showstringspaces=false
}
\lstdefinelanguage{typescript}{
  morekeywords={import,export,from,const,let,var,function,return,if,else,for,of,in,while,switch,case,break,async,await,try,catch,finally,throw,new,class,this,extends,implements,interface,type,enum,readonly,public,private,protected,static,get,set,true,false,null,undefined,void,string,number,boolean,any,unknown,never,typeof,keyof,as},
  sensitive=true,
  morecomment=[l]{//},
  morecomment=[s]{/*}{*/},
  morestring=[b]',
  morestring=[b]",
  morestring=[b]`
}
\lstalias{javascript}{typescript}
\lstset{
  backgroundcolor=\color{codebg},
  basicstyle=\ttfamily\footnotesize,
  frame=single,
  rulecolor=\color{codeframeclr},
  frameround=tttt,
  breaklines=true,
  tabsize=2,
  showstringspaces=false,
  upquote=true,
  columns=fullflexible
}

% ---------- Helper macros ----------
\newcommand{\codefile}[1]{\texttt{#1}}
\newcommand{\kbd}[1]{\texttt{#1}}

% ---------- Code frame environment ----------
% Automatski dodaje fragile i shrink=5, koristi se umjesto \begin{frame}
% Primjer: \begin{codeframe}{Naslov} ... \end{codeframe}
% Primjer s opcijama: \begin{codeframe}[shrink=15]{Naslov} ... \end{codeframe}
\newenvironment{codeframe}[2][]
  {\begin{frame}[fragile,shrink=5,#1]{#2}}
  {\end{frame}}

% ============================================================

\documentclass[aspectratio=149,11pt]{beamer}
\usetheme{CambridgeUS}

\usepackage{fontspec}              % radi samo sa XeLaTeX/LuaLaTeX
\setsansfont{Latin Modern Sans}    % sans-serif font
\setmonofont{Latin Modern Mono}    % monospace font za kod
\usepackage[croatian]{babel}
\usepackage{csquotes}              % \enquote{} → hrvatski navodnici „..."
\usepackage{lmodern}
\usepackage{booktabs}
\usepackage{listings}
\usepackage{xcolor}
\usepackage{hyperref}

% ---------- Autor i logo ----------
\author{mag.ing. Josip Torić (SIT UNIZD)}
\titlegraphic{\includegraphics[height=2.5cm]{../common/SIT_logo_tekst1.pdf}}

% ---------- Listings ----------
\definecolor{codebg}{RGB}{248,248,248}
\definecolor{codeframeclr}{RGB}{220,220,220}
\lstdefinelanguage{bash}{
  morekeywords={sudo,apt,brew,pip,pipx,uv,winget,python,git,flask,source,deactivate,code,gunicorn,docker,compose,heroku,render,railway,node,npm,pnpm,yarn,uvicorn,alembic,pytest,ruff,black,precommit},
  sensitive=false,
  morecomment=[l]{\#}
}
\lstdefinelanguage{python}{
  morekeywords={from,import,def,class,return,if,elif,else,for,in,while,try,except,raise,with,as,not,and,or,is,lambda,yield,True,False,None,await,async},
  sensitive=true,
  morecomment=[l]{\#},
  morestring=[b]',
  morestring=[b]"
}
\lstdefinelanguage{sql}{
  morekeywords={SELECT,FROM,WHERE,INSERT,INTO,VALUES,UPDATE,SET,DELETE,CREATE,TABLE,PRIMARY,KEY,FOREIGN,REFERENCES,UNIQUE,NOT,NULL,INDEX,ON,ALTER,DROP,CONSTRAINT},
  sensitive=false,
  morecomment=[l]{--}
}
\lstdefinelanguage{json}{
  morecomment=[l]{//},
  morestring=[b]",
  showstringspaces=false
}
\lstdefinelanguage{typescript}{
  morekeywords={import,export,from,const,let,var,function,return,if,else,for,of,in,while,switch,case,break,async,await,try,catch,finally,throw,new,class,this,extends,implements,interface,type,enum,readonly,public,private,protected,static,get,set,true,false,null,undefined,void,string,number,boolean,any,unknown,never,typeof,keyof,as},
  sensitive=true,
  morecomment=[l]{//},
  morecomment=[s]{/*}{*/},
  morestring=[b]',
  morestring=[b]",
  morestring=[b]`
}
\lstalias{javascript}{typescript}
\lstset{
  backgroundcolor=\color{codebg},
  basicstyle=\ttfamily\footnotesize,
  frame=single,
  rulecolor=\color{codeframeclr},
  frameround=tttt,
  breaklines=true,
  tabsize=2,
  showstringspaces=false,
  upquote=true,
  columns=fullflexible
}

% ---------- Helper macros ----------
\newcommand{\codefile}[1]{\texttt{#1}}
\newcommand{\kbd}[1]{\texttt{#1}}

% ---------- Code frame environment ----------
% Automatski dodaje fragile i shrink=5, koristi se umjesto \begin{frame}
% Primjer: \begin{codeframe}{Naslov} ... \end{codeframe}
% Primjer s opcijama: \begin{codeframe}[shrink=15]{Naslov} ... \end{codeframe}
\newenvironment{codeframe}[2][]
  {\begin{frame}[fragile,shrink=5,#1]{#2}}
  {\end{frame}}

% ============================================================

\documentclass[aspectratio=149,11pt]{beamer}
\usetheme{CambridgeUS}

\usepackage{fontspec}              % radi samo sa XeLaTeX/LuaLaTeX
\setsansfont{Latin Modern Sans}    % sans-serif font
\setmonofont{Latin Modern Mono}    % monospace font za kod
\usepackage[croatian]{babel}
\usepackage{csquotes}              % \enquote{} → hrvatski navodnici „..."
\usepackage{lmodern}
\usepackage{booktabs}
\usepackage{listings}
\usepackage{xcolor}
\usepackage{hyperref}

% ---------- Autor i logo ----------
\author{mag.ing. Josip Torić (SIT UNIZD)}
\titlegraphic{\includegraphics[height=2.5cm]{../common/SIT_logo_tekst1.pdf}}

% ---------- Listings ----------
\definecolor{codebg}{RGB}{248,248,248}
\definecolor{codeframeclr}{RGB}{220,220,220}
\lstdefinelanguage{bash}{
  morekeywords={sudo,apt,brew,pip,pipx,uv,winget,python,git,flask,source,deactivate,code,gunicorn,docker,compose,heroku,render,railway,node,npm,pnpm,yarn,uvicorn,alembic,pytest,ruff,black,precommit},
  sensitive=false,
  morecomment=[l]{\#}
}
\lstdefinelanguage{python}{
  morekeywords={from,import,def,class,return,if,elif,else,for,in,while,try,except,raise,with,as,not,and,or,is,lambda,yield,True,False,None,await,async},
  sensitive=true,
  morecomment=[l]{\#},
  morestring=[b]',
  morestring=[b]"
}
\lstdefinelanguage{sql}{
  morekeywords={SELECT,FROM,WHERE,INSERT,INTO,VALUES,UPDATE,SET,DELETE,CREATE,TABLE,PRIMARY,KEY,FOREIGN,REFERENCES,UNIQUE,NOT,NULL,INDEX,ON,ALTER,DROP,CONSTRAINT},
  sensitive=false,
  morecomment=[l]{--}
}
\lstdefinelanguage{json}{
  morecomment=[l]{//},
  morestring=[b]",
  showstringspaces=false
}
\lstdefinelanguage{typescript}{
  morekeywords={import,export,from,const,let,var,function,return,if,else,for,of,in,while,switch,case,break,async,await,try,catch,finally,throw,new,class,this,extends,implements,interface,type,enum,readonly,public,private,protected,static,get,set,true,false,null,undefined,void,string,number,boolean,any,unknown,never,typeof,keyof,as},
  sensitive=true,
  morecomment=[l]{//},
  morecomment=[s]{/*}{*/},
  morestring=[b]',
  morestring=[b]",
  morestring=[b]`
}
\lstalias{javascript}{typescript}
\lstset{
  backgroundcolor=\color{codebg},
  basicstyle=\ttfamily\footnotesize,
  frame=single,
  rulecolor=\color{codeframeclr},
  frameround=tttt,
  breaklines=true,
  tabsize=2,
  showstringspaces=false,
  upquote=true,
  columns=fullflexible
}

% ---------- Helper macros ----------
\newcommand{\codefile}[1]{\texttt{#1}}
\newcommand{\kbd}[1]{\texttt{#1}}

% ---------- Code frame environment ----------
% Automatski dodaje fragile i shrink=5, koristi se umjesto \begin{frame}
% Primjer: \begin{codeframe}{Naslov} ... \end{codeframe}
% Primjer s opcijama: \begin{codeframe}[shrink=15]{Naslov} ... \end{codeframe}
\newenvironment{codeframe}[2][]
  {\begin{frame}[fragile,shrink=5,#1]{#2}}
  {\end{frame}}


\title[Razvoj web aplikacija]{Razvoj web aplikacija - Predavanje 1 \\ Backend arhitektura i skeleton projekta}
\date{05.10.2026.}

\begin{document}

\frame{\titlepage}

% ------------------------------------------------------------
\begin{frame}{Plan predavanja}
\begin{enumerate}
  \item Što znači \enquote{backend arhitektura} (i što je \emph{loša} arhitektura)
  \item Kako izgleda tipičan request: frontend $\rightarrow$ API $\rightarrow$ DB
  \item Skeleton projekta: struktura repozitorija i moduli
  \item Dev okruženje: Docker Compose, env varovi, konfiguracija
  \item FastAPI: app factory, routeri, dependencies, health endpoint
  \item Git workflow: commitovi, tagovi i \texttt{git diff} kao alat za učenje
\end{enumerate}

\vspace{0.6em}
Na kraju predavanja imate \textbf{pokrenut} projekt: baza + API + dokumentiran način pokretanja.
\end{frame}

% ------------------------------------------------------------
\begin{frame}{Ishodi učenja (što morate znati objasniti)}
Nakon ovog predavanja student može:
\begin{itemize}
  \item objasniti tok HTTP zahtjeva kroz slojeve backend aplikacije
  \item razlikovati \textbf{router}, \textbf{service} i \textbf{DB layer}
  \item postaviti razvojno okruženje: Docker Compose + env varovi
  \item objasniti zašto koristimo \textbf{app factory} i modularnu strukturu
  \item koristiti \texttt{git diff} za razumijevanje promjena u kodu
\end{itemize}
\end{frame}

% ------------------------------------------------------------
\begin{frame}{Zašto se uopće bavimo arhitekturom?}
\begin{itemize}
  \item Jer \enquote{radi mi} nije isto što i \enquote{održivo je}
  \item Jer na obrani ispitujem \textbf{razumijevanje}, ne copy/paste
  \item Jer se problemi pojavljuju kad projekt naraste:
  \begin{itemize}
    \item više endpointa, više pravila, više korisnika
    \item više developera (vi u timu)
    \item više bugova i edge caseova
  \end{itemize}
\end{itemize}

\vspace{0.6em}
\textbf{Arhitektura} = kako je sustav organiziran i gdje živi odgovornost.
\end{frame}

% ------------------------------------------------------------
\begin{frame}{Najčešći anti-pattern: \enquote{Sve u routeru}}
\begin{columns}[T]
\column{0.48\textwidth}
\textbf{Sve u routeru}
\begin{itemize}
  \item jedna velika datoteka: HTTP + validacija + DB upiti + poslovna pravila + errori
  \item ne možeš testirati bez HTTP-a
  \item svaka promjena razbije nešto drugo
\end{itemize}

\column{0.48\textwidth}
\textbf{Slojevi}
\begin{itemize}
  \item router je tanak: samo HTTP
  \item poslovna logika u service-u
  \item DB detalji u repository/DB layeru
\end{itemize}
\end{columns}

\vspace{0.8em}
Na obrani: tko ima sve u routeru, ne zna objasniti \emph{što je gdje} i \emph{zašto}.
\end{frame}

% ------------------------------------------------------------
\begin{frame}{Request flow (mentalni model)}
\begin{enumerate}
  \item Frontend šalje HTTP request (npr. \texttt{POST /auth/login})
  \item Router parsira request + provjerava auth dependency
  \item Service provodi poslovna pravila (npr. rokovi, ownership)
  \item Repository/DB layer radi upit (transakcija)
  \item Service vraća rezultat (ili baca domain error)
  \item Router mapira rezultat u HTTP response (status code + JSON)
\end{enumerate}

\vspace{0.6em}
HTTP (router) je samo \textbf{transport}. Business pravila su \textbf{domena}.
\end{frame}

% ------------------------------------------------------------
\begin{frame}{REST u praksi (minimum koji tražim)}
Resource-oriented URL-ovi: \texttt{GET /lifters}, \texttt{POST /lifters}, \texttt{GET /lifters/\{id\}}, \texttt{PATCH /lifters/\{id\}}, \texttt{DELETE /lifters/\{id\}}.

\vspace{0.6em}
Status codeovi nisu dekoracija:

\vspace{0.4em}
\begin{tabular}{@{}lp{7.5cm}@{}}
  \toprule
  \textbf{Kod} & \textbf{Kada} \\
  \midrule
  200 / 201 & uspjeh / resurs kreiran \\
  400 / 422 & neispravan zahtjev / validacija inputa pala \\
  401 / 403 & nisi prijavljen / nemaš ovlasti \\
  404 & resurs ne postoji \\
  409 & konflikt, npr. duplikat prijave \\
  500 & greška servera \\
  \bottomrule
\end{tabular}

\vspace{0.6em}
\texttt{GET}, \texttt{PUT} i \texttt{DELETE} su u pravilu idempotentni.
\end{frame}

% ------------------------------------------------------------
\begin{frame}{DTO i validacija (zašto Pydantic)}
Backend mora validirati input - frontend se ne računa kao zaštita. DTO, odnosno schema, je ugovor API-ja: definira što točno ulazi i što izlazi.

\vspace{0.6em}
Pydantic taj ugovor provodi automatski: nevaljan input dobije jasan 422 s opisom greške, a iz istih schema se generiraju OpenAPI modeli.

\vspace{0.8em}
\begin{block}{Na obrani pitam}
Gdje je schema, što validira, koji error vraća i zašto.
\end{block}
\end{frame}

% ------------------------------------------------------------
\begin{frame}{Konfiguracija: 12-factor i env varovi}
\begin{itemize}
  \item Konfiguracija ne smije biti hardkodirana
  \item Tajne ne idu u git:
  \begin{itemize}
    \item DB password, JWT secret, API base URL
  \end{itemize}
  \item \codefile{.env} je za lokalno, \codefile{.env.example} je za repo
  \item U produkciji: env varovi dolaze iz platforme (Railway/Render)
\end{itemize}

\vspace{0.6em}
Ako promijeniš okruženje (lokalno $\rightarrow$ produkcija), ne smiješ mijenjati kod.
\end{frame}

% ------------------------------------------------------------
\begin{frame}{Docker Compose: zašto ga koristimo?}
Standardizira PostgreSQL za sve studente, uklanja \enquote{meni ne radi na Windowsu/Macu} i omogućava brzo resetiranje okruženja.

\vspace{0.8em}
\begin{tabular}{@{}p{5.2cm}p{5.2cm}@{}}
  \toprule
  \textbf{U Composeu} & \textbf{Lokalno} \\
  \midrule
  \texttt{db} - Postgres + volume & backend - \texttt{uvicorn --reload} \\
   & frontend - Vite dev server \\
  \bottomrule
\end{tabular}

\vspace{0.6em}
Backend i frontend ne dockeriziramo radi brzine razvoja.
\end{frame}

% ------------------------------------------------------------
\begin{frame}[fragile,shrink=5]{Repo struktura (monorepo)}
\begin{lstlisting}[language=bash]
repo/
  api/              # FastAPI backend
  web/              # Vue frontend
  docker-compose.yml
  .gitignore
  README.md
  .env.example
\end{lstlisting}

\vspace{0.6em}
Jedan repo = lakša integracija, jedan README, jedna priča na obrani.
\end{frame}

% ------------------------------------------------------------
\begin{frame}[fragile,shrink=15]{Backend struktura (minimalno, ali ozbiljno)}
\begin{lstlisting}[language=bash]
api/
  app/
    main.py            # app factory + include routers
    core/
      config.py        # env config
      deps.py          # dependencies (auth, db)
      errors.py        # custom errors / handlers
      logging.py       # logging config
    routers/
      health.py
    schemas/
    models/
    services/
    repositories/
  tests/
  pyproject.toml (ili requirements.txt)
  alembic/ (kasnije)
\end{lstlisting}
\end{frame}

% ------------------------------------------------------------
\begin{frame}{App factory: zašto?}
\begin{itemize}
  \item Omogućava:
  \begin{itemize}
    \item testiranje bez pokretanja servera
    \item konfiguraciju ovisno o okruženju
    \item modularno dodavanje routera/middlewarea
  \end{itemize}
  \item Smanjuje \enquote{global state} probleme
\end{itemize}

\vspace{0.6em}
\texttt{create\_app()} napravi aplikaciju. \texttt{uvicorn} samo \textbf{servira}.
\end{frame}

% ------------------------------------------------------------
\begin{frame}[fragile,shrink=5]{Primjer: minimalni \texttt{main.py}}
\begin{lstlisting}[language=python]
from fastapi import FastAPI
from app.routers.health import router as health_router

def create_app() -> FastAPI:
    app = FastAPI(title="Powerlifting Registrations API")
    app.include_router(health_router, prefix="/health", tags=["health"])
    return app

app = create_app()
\end{lstlisting}
\end{frame}

% ------------------------------------------------------------
\begin{frame}{Health endpoint (zašto je bitan?)}
\begin{itemize}
  \item Najjednostavniji endpoint koji potvrđuje:
  \begin{itemize}
    \item da se aplikacija diže
    \item da routing radi
    \item (kasnije) da DB konekcija radi
  \end{itemize}
  \item U produkciji se često koristi kao readiness/liveness check
\end{itemize}

\vspace{0.6em}
\texttt{GET /health} $\rightarrow$ \texttt{\{ "status": "ok" \}}.
\end{frame}

% ------------------------------------------------------------
\begin{frame}[fragile,shrink=5]{Primjer: \texttt{routers/health.py}}
\begin{lstlisting}[language=python]
from fastapi import APIRouter

router = APIRouter()

@router.get("/")
def health():
    return {"status": "ok"}
\end{lstlisting}
\end{frame}

% ------------------------------------------------------------
\begin{frame}[fragile,shrink=5]{Konfiguracija: \texttt{core/config.py}}
\begin{itemize}
  \item Čitamo env varove na jednom mjestu
  \item Ne razbacujemo \texttt{os.getenv} po kodu
\end{itemize}

\begin{lstlisting}[language=python]
from pydantic_settings import BaseSettings

class Settings(BaseSettings):
    ENV: str = "dev"
    DATABASE_URL: str
    JWT_SECRET: str = "change-me"
    JWT_ISSUER: str = "sit-unizd"

settings = Settings()
\end{lstlisting}

\vspace{0.4em}
Pydantic settings je najčišći put za konfiguraciju.
\end{frame}

% ------------------------------------------------------------
\begin{frame}{Logging (što logiramo, što ne?)}
\begin{tabular}{@{}p{5.2cm}p{5.2cm}@{}}
  \toprule
  \textbf{Logiramo} & \textbf{Ne logiramo} \\
  \midrule
  request metadata: ruta, status code & lozinke, tokene, osobne identifikatore \\
  \midrule
  error stack trace u dev okruženju & kompletan request body ako sadrži osjetljive podatke \\
  \midrule
  ključne domenske događaje, npr. \enquote{registration withdrawn} & \\
  \bottomrule
\end{tabular}

\vspace{0.6em}
Log nije baza. Log je trag za debugging.
\end{frame}

% ------------------------------------------------------------
\begin{frame}[fragile,shrink=5]{Logging: minimalna konfiguracija}
\begin{lstlisting}[language=python]
# app/core/logging.py
import logging

def setup_logging(debug: bool = False) -> None:
    logging.basicConfig(
        level=logging.DEBUG if debug else logging.INFO,
        format="%(asctime)s %(levelname)s %(name)s: %(message)s",
    )

# u servisu:
logger = logging.getLogger(__name__)
logger.info("club created: id=%s", club.id)
\end{lstlisting}

\vspace{0.4em}
Logger se dohvaća po modulu preko \texttt{\_\_name\_\_} pa svaki zapis kaže odakle dolazi. U produkciji sve ide na stdout, a platforma skuplja logove - detalji u Predavanju 12.
\end{frame}

% ------------------------------------------------------------
\begin{frame}{Error handling (na razini arhitekture)}
\begin{itemize}
  \item Dosljedan format errora:
  \begin{itemize}
    \item \texttt{code}, \texttt{message}, (opc.) \texttt{fields}
  \end{itemize}
  \item Razdvajamo:
  \begin{itemize}
    \item \textbf{domain error} (poslovno pravilo) vs \textbf{HTTP error}
  \end{itemize}
  \item Kasnije: global exception handler mapira domain error u status code
\end{itemize}

\vspace{0.6em}
Ako to ne dogovorite odmah, poslije je kaos.
\end{frame}

% ------------------------------------------------------------
\begin{frame}{Tooling: lint i format (zašto Ruff?)}
\begin{itemize}
  \item \textbf{Ruff} je danas i linter i formatter, jedan alat za oboje:
  \begin{itemize}
    \item \texttt{ruff check .} - lint: importi, greške, stil
    \item \texttt{ruff format .} - formatiranje, jedan stil bez rasprave
  \end{itemize}
  \item Formatter je naslijedio Black i drži isti stil koda
  \item Smanjuje code review overhead i olakšava timski rad
\end{itemize}

\vspace{0.6em}
Ako u repou vidim nered, na obrani se prvo bavimo time i gubite vrijeme.
\end{frame}

% ------------------------------------------------------------
\begin{frame}{Testiranje: što znači \enquote{testabilna arhitektura}?}
\begin{itemize}
  \item Ako je sva logika u routeru, testiraš sve preko HTTP-a $\Rightarrow$ sporo i krhko
  \item Ako je logika u service-u, možeš testirati:
  \begin{itemize}
    \item pravila (rokovi, ownership) bez HTTP sloja
  \end{itemize}
  \item FastAPI omogućava i integration testove (TestClient/httpx)
\end{itemize}

\vspace{0.6em}
Danas postavljamo skeleton tako da se testovi kasnije pišu bez muke.
\end{frame}

% ------------------------------------------------------------
\begin{frame}{Git kao nastavna metoda: commitovi i diff}
\begin{center}
\Large Svaki commit je mala lekcija. \\[0.3em] Veliki commit = loša lekcija.
\end{center}

\vspace{0.8em}
Ne live-codeam: pokazujem unaprijed pripremljen projekt kroz commitove. Vi učite gledajući \textbf{što se promijenilo} i \textbf{zašto} - \texttt{git diff} je alat za razumijevanje, ne samo za rješavanje konflikata.
\end{frame}

% ------------------------------------------------------------
\begin{frame}{Kako čitati diff (brza pravila)}
\begin{itemize}
  \item Prvo razumij \textbf{što} je cilj promjene (commit message)
  \item Onda gledaj:
  \begin{itemize}
    \item novi fileovi vs izmjene postojećih
    \item javne API promjene (npr. novi endpoint)
    \item promjene u configu (env, compose)
  \end{itemize}
  \item Tek na kraju detalji implementacije
\end{itemize}
\end{frame}

% ------------------------------------------------------------
\begin{frame}[fragile,shrink=5]{Korisne git naredbe}
\begin{lstlisting}[language=bash]
# pregled povijesti
git log --oneline --decorate

# diff izmedu commitova
git diff <commit1>..<commit2>

# diff u konkretnom commitu
git show <commit>

# promjena na konkretan tag
git checkout lecture-01-end
\end{lstlisting}
\end{frame}

% ------------------------------------------------------------
\begin{frame}{Commitovi za Predavanje 1 (što obrađujemo danas)}
\begin{itemize}
  \item \texttt{chore: init monorepo (api, web) with tooling configs}
  \item \texttt{chore: add docker-compose postgres and env example files}
  \item \texttt{feat(api): bootstrap FastAPI app with health check and version}
  \item \texttt{refactor(api): add app factory and basic module structure}
  \item \texttt{docs: add local dev guide and project conventions}
\end{itemize}

\vspace{0.6em}
\textbf{Na satu:} prolazimo diffeve i objašnjavamo \emph{zašto} je to u toj datoteci.
\end{frame}

% ------------------------------------------------------------
\begin{frame}[fragile,shrink=5]{Demo: \texttt{docker-compose.yml} (minimalno)}
\begin{lstlisting}[language=bash]
services:
  db:
    image: postgres:17
    environment:
      POSTGRES_DB: pl_reg
      POSTGRES_USER: pl_user
      POSTGRES_PASSWORD: pl_pass
    ports:
      - "5432:5432"
    volumes:
      - pgdata:/var/lib/postgresql/data

volumes:
  pgdata:
\end{lstlisting}

\vspace{0.4em}
Lozinke su demo - u praksi se stavljaju u \codefile{.env}.
\end{frame}

% ------------------------------------------------------------
\begin{frame}[fragile,shrink=5]{Pokretanje baze i API-ja (danas)}
\begin{lstlisting}[language=bash]
# baza
docker compose up -d db

# backend
cd api
uv venv
source .venv/bin/activate
uv pip install -r requirements.txt
uvicorn app.main:app --reload
\end{lstlisting}

\vspace{0.4em}
\texttt{uv} objedinjuje venv i pip, brži je i danas je standard za nove projekte - \codefile{requirements.txt} ostaje isti.

\vspace{0.4em}
Checkpoint: otvorite \texttt{http://127.0.0.1:8000/health} i dobijte \texttt{"ok"}.
\end{frame}

% ------------------------------------------------------------
\begin{frame}{Swagger / OpenAPI (zašto je bitno od prvog dana?)}
\begin{itemize}
  \item API dokumentacija se generira automatski
  \item Olakšava testiranje endpointa
  \item Frontend se oslanja na ugovor API-ja
\end{itemize}

\vspace{0.6em}
\texttt{/docs} i \texttt{/openapi.json} su vaši prijatelji.
\end{frame}

% ------------------------------------------------------------
\begin{frame}{Što ide sljedeće (Predavanje 2)}
\begin{itemize}
  \item Modeliranje podataka: Club i User
  \item SQLAlchemy engine + session
  \item Alembic migracije (inicijalno)
  \item Seed admina i demo kluba
\end{itemize}

\vspace{0.8em}
\begin{center}
\Large Pitanja?
\end{center}
\end{frame}

\end{document}
